\documentclass[journal]{IEEEtran}

% ============================================================
% Packages
% ============================================================
\usepackage{amsmath,amssymb,amsfonts}
\usepackage{algorithmic}
\usepackage{algorithm}
\usepackage{graphicx}
\usepackage{textcomp}
\usepackage{xcolor}
\usepackage{booktabs}
\usepackage{multirow}
\usepackage{hyperref}
\usepackage{subcaption}
\usepackage{balance}
\usepackage{tabularx}
\usepackage{array}

% For placeholder results
\newcommand{\TODO}[1]{\textcolor{red}{\textbf{[#1]}}}

\begin{document}

% ============================================================
% Title
% ============================================================
\title{Modeling Cascade Contagion in Decentralized Finance as a Computational Social System via Temporal Graph Networks}

\author{%
\IEEEauthorblockN{Keshav Krishnan}
% \IEEEauthorblockA{Department \\ University \\ email@university.edu}
}

\maketitle

% ============================================================
% Abstract
% ============================================================
\begin{abstract}
Decentralized Finance (DeFi) constitutes a computational social system in which autonomous protocol communities---governed through decentralized autonomous organizations (DAOs)---collectively allocate capital, manage risk, and respond to crises through interconnected smart contract infrastructure.
When stress events occur, these sociotechnical interdependencies can trigger liquidation cascades: self-reinforcing contagion processes that propagate across protocols through shared collateral, liquidity pathways, oracle dependencies, and governance overlaps, causing billions of dollars in losses.
We propose a \emph{Temporal Graph Network} (TGN) framework that models the DeFi ecosystem as a heterogeneous composability graph with six distinct edge types capturing cross-protocol dependencies.
Our model combines per-node GRU memory for tracking protocol state evolution, heterogeneous graph attention for modeling contagion pathways, and temporal feature augmentation for capturing pre-cascade warning signals.
We formulate cascade prediction as a multi-horizon binary classification task at 1-day, 3-day, 7-day, and 30-day horizons, enabling both operational alerting and strategic risk planning.
Evaluated on real-world Total Value Locked (TVL) data spanning 1,734 days across 15 major Ethereum DeFi protocols---encompassing six verified historical cascade events and seven statistically detected anomalies---our TGN is compared against five baselines: Static GNN, LSTM, XGBoost, SIR contagion, and network centrality.
Ablation studies quantify the complementary contributions of graph topology, temporal memory, and multi-scale feature engineering.
Statistical validation via Diebold-Mariano tests and bootstrap confidence intervals establishes significance.
Our framework provides an actionable early warning system for DeFi risk management and contributes a principled methodology for analyzing cascade contagion dynamics in computational social systems.
\end{abstract}

\begin{IEEEkeywords}
Computational social systems, DeFi, liquidation cascades, temporal graph networks, systemic risk, composability graph, contagion modeling, decentralized governance
\end{IEEEkeywords}

% ============================================================
% I. Introduction
% ============================================================
\section{Introduction}

\IEEEPARstart{D}{ecentralized} Finance (DeFi) has grown into a complex financial ecosystem with over \$100 billion in Total Value Locked (TVL) across hundreds of interconnected protocols~\cite{werner2022defi}.
Unlike traditional finance, where systemic risk is mediated through regulated clearinghouses and central counterparties, DeFi protocols interact through permissionless smart contracts that create deeply entangled dependency structures.
Protocols share collateral assets (e.g., WETH, WBTC, USDC), route liquidity through decentralized exchanges, depend on common oracle price feeds, and overlap in governance token holdings.

These interdependencies create \emph{composability}---the ability for protocols to build upon each other---which is simultaneously DeFi's greatest innovation and its primary source of systemic risk.
When a stress event occurs at one protocol, the consequences propagate through these dependency channels, potentially triggering a \emph{liquidation cascade}: a self-reinforcing cycle of collateral liquidations, price drops, and further liquidations that spreads across the ecosystem.

Crucially, DeFi constitutes a \emph{computational social system}~\cite{wang2020computational} in which human communities govern shared financial infrastructure through decentralized autonomous organizations (DAOs), collective capital allocation decisions, and trust-mediated protocol interactions.
Cascade events are not merely technical failures but \emph{social contagion} processes: panic selling driven by herd behavior~\cite{da2022crypto}, governance failures that amplify rather than contain crises, and collective loss of confidence that propagates through community networks.
The composability graph that connects DeFi protocols is fundamentally a \emph{social} artifact---designed by human communities, reflecting their collective risk preferences, and shaped by the governance decisions of token holders.
Understanding cascade contagion in DeFi thus requires computational methods that jointly model the social governance structure and the technical interdependency network.

Historical events demonstrate the severity of this risk.
The Terra/Luna collapse in May 2022 erased over \$40 billion in value within days, with contagion spreading to lending protocols (Aave, Compound), stablecoins (MakerDAO), and centralized entities (Three Arrows Capital, Celsius)~\cite{briola2023anatomy}.
The FTX collapse in November 2022 similarly cascaded through DeFi, causing 25\% TVL drops across major protocols.
More recently, the Curve exploit in July 2023 demonstrated that even moderate-severity events can trigger disproportionate systemic effects through governance token contagion.

Despite the critical importance of cascade prediction, existing approaches suffer from three key limitations:
\begin{enumerate}
    \item \textbf{Graph structure neglect.} Most risk models treat protocols independently, ignoring the composability graph that defines contagion pathways~\cite{zhou2023sok}.
    \item \textbf{Static analysis.} Network-based approaches typically analyze graph topology at a single point in time, failing to capture the temporal evolution of protocol states that precede cascades~\cite{qin2021cefi}.
    \item \textbf{Single-horizon prediction.} Existing models predict cascades at a single time horizon, whereas practitioners require multi-scale predictions for different operational decisions~\cite{xu2023defi}.
\end{enumerate}

We address these limitations with a \emph{Temporal Graph Network} (TGN) framework that jointly models the DeFi composability graph structure and the temporal dynamics of protocol state evolution.
Our contributions are:

\begin{enumerate}
    \item \textbf{DeFi Composability Graph.} We formalize the cross-protocol dependency structure as a heterogeneous graph with six semantically distinct edge types: shared collateral, liquidity flow, oracle dependency, governance overlap, price correlation, and liquidation pathways. This graph captures the multi-channel nature of DeFi contagion as a computational social system.

    \item \textbf{Temporal Graph Network for Cascade Prediction.} We propose a TGN architecture that combines per-node GRU memory, heterogeneous graph attention, temporal feature augmentation, and multi-horizon prediction heads. The model processes daily snapshots of 15 major DeFi protocols, each described by 46 features spanning TVL dynamics, price movements, lending markets, network topology, and macroeconomic indicators.

    \item \textbf{Real-World Evaluation.} We evaluate on real TVL data from DeFiLlama spanning 1,734 days (June 2021--February 2026), encompassing six verified historical cascade events and additional statistically detected anomalies. We formulate prediction at four horizons (1-day, 3-day, 7-day, 30-day) with rigorous temporal splitting to prevent data leakage.

    \item \textbf{Comprehensive Benchmarking.} We compare against five baselines representing graph-only (Static GNN), temporal-only (LSTM), traditional ML (XGBoost), physics-based (SIR contagion), and network-analytic (centrality) approaches, with rigorous statistical testing and ablation studies.
\end{enumerate}

% ============================================================
% II. Related Work
% ============================================================
\section{Related Work}

\subsection{Computational Social Systems and Decentralized Finance}

Computational social systems~\cite{wang2020computational} encompass man-machine interactive systems where human social behavior and computational processes are deeply intertwined.
DeFi represents a paradigmatic computational social system: protocols are governed by DAOs where token holders vote on risk parameters, capital allocation is determined by collective human decisions mediated by smart contracts, and systemic risk emerges from the aggregate behavior of millions of interacting agents.

Research on social dynamics in cryptocurrency markets has demonstrated the prevalence of herd behavior~\cite{da2022crypto}, information cascades in decentralized governance~\cite{fritsch2022dao}, and the role of social sentiment in driving price dynamics~\cite{chen2023social}.
Ante~\cite{ante2023defi} analyzes how DeFi protocol interactions create systemic risk through social and technical channels.
Li et al.~\cite{li2023cascade} study cascade dynamics in online social networks, providing methodological foundations that parallel our DeFi contagion analysis.
Wang et al.~\cite{wang2022social} examine how social network structure influences cascade propagation, directly relevant to DeFi where governance communities overlap across protocols.

Our work positions DeFi cascade prediction as a computational social systems problem, arguing that the composability graph connecting protocols is a social construct reflecting collective governance decisions, and that cascade propagation is mediated by both technical dependencies and social contagion dynamics.

\subsection{Systemic Risk in DeFi}

The study of systemic risk in DeFi has attracted growing attention as the ecosystem matures.
Werner et al.~\cite{werner2022defi} provide a foundational survey of DeFi's structure and risks.
Gudgeon et al.~\cite{gudgeon2020defi} analyze the vulnerability of DeFi protocols to governance attacks and economic exploits.
Qin et al.~\cite{qin2021cefi} study the interaction between centralized and decentralized finance during stress events.
Briola et al.~\cite{briola2023anatomy} provide a detailed anatomy of the Terra/Luna crash and its contagion effects, demonstrating how interconnected DeFi positions amplify shocks.
Zhou et al.~\cite{zhou2023sok} systematize knowledge on DeFi attacks and identify composability as a key risk amplifier.

Several works have studied specific contagion channels.
Perez et al.~\cite{perez2021liquidations} analyze liquidation mechanisms in lending protocols and their feedback effects on asset prices.
Xu et al.~\cite{xu2023defi} study cross-protocol dependencies through shared collateral and identify systemic risk concentration.
Aramonte et al.~\cite{aramonte2022defi} highlight the parallels between DeFi's interconnectedness and traditional financial systemic risk.

\subsection{Graph Neural Networks for Financial Risk}

Graph neural networks (GNNs) have been successfully applied to financial risk modeling.
Wang et al.~\cite{wang2021review} survey GNN applications in finance, including fraud detection, credit scoring, and portfolio optimization.
Cheng et al.~\cite{cheng2022financial} apply GNNs to financial network analysis, demonstrating that graph structure captures contagion risk better than node features alone.
Yang et al.~\cite{yang2020financial} use heterogeneous GNNs for financial distress prediction, showing that multi-relational edges improve over homogeneous graphs.

In the specific context of blockchain and DeFi, Li et al.~\cite{li2022tgn} apply temporal graph networks to Ethereum transaction analysis.
Weber et al.~\cite{weber2019anti} use GNNs on Bitcoin transaction graphs for anti-money laundering.
Kanezashi et al.~\cite{kanezashi2022ethereum} study Ethereum account classification using graph attention networks.

\subsection{Temporal Graph Networks}

Temporal graph networks, introduced by Rossi et al.~\cite{rossi2020tgn}, extend static GNNs by incorporating time-stamped interactions and per-node memory modules.
The key innovation is a memory module that maintains a compressed representation of each node's history, updated through message passing at each interaction.
This architecture has shown strong performance on dynamic link prediction and node classification tasks.

Xu et al.~\cite{xu2020inductive} propose TGAT (Temporal Graph Attention), which uses self-attention over temporal neighborhoods.
Kazemi et al.~\cite{kazemi2020representation} survey representation learning for dynamic graphs.
Pareja et al.~\cite{pareja2020evolvegcn} propose EvolveGCN, which uses an RNN to evolve GCN parameters over time.

Our work differs from prior TGN applications in three ways: (1) we model a \emph{heterogeneous} composability graph with six semantically distinct edge types rather than a single interaction type; (2) we incorporate domain-specific temporal feature augmentation; and (3) we target a multi-horizon prediction task specific to financial cascade events in a computational social system.

\subsection{Contagion and Cascade Models}

Epidemiological models, particularly the SIR (Susceptible-Infected-Recovered) framework, have been adapted for financial contagion.
Hurd et al.~\cite{hurd2016contagion} apply SIR-type models to interbank networks.
Amini et al.~\cite{amini2016resilience} analyze the resilience of financial networks to cascading failures using percolation theory.
In DeFi, contagion models have been applied by Tolmach et al.~\cite{tolmach2021formal} who formally verify DeFi protocol interactions, and by Bartoletti et al.~\cite{bartoletti2021formal} who model composability risks.

Our SIR baseline implements a standard compartmental model on the composability graph adjacency matrix, providing a physics-based comparison for our data-driven TGN approach.


% ============================================================
% III. Problem Formulation
% ============================================================
\section{Problem Formulation}

\subsection{DeFi Ecosystem as a Temporal Graph}

We model the DeFi ecosystem as a temporal heterogeneous graph $\mathcal{G}_t = (\mathcal{V}, \mathcal{E}, \mathbf{X}_t)$ at each discrete timestep $t \in \{1, \ldots, T\}$, where:
\begin{itemize}
    \item $\mathcal{V} = \{v_1, \ldots, v_N\}$ is the set of $N = 15$ DeFi protocol nodes.
    \item $\mathcal{E} = \{e_1, \ldots, e_K\}$ is a set of $K = 6$ heterogeneous edge types representing distinct cross-protocol dependencies.
    \item $\mathbf{X}_t \in \mathbb{R}^{N \times F}$ is the node feature matrix at time $t$, where $F = 46$ features describe each protocol's state.
\end{itemize}

The edge structure is static (reflecting persistent architectural dependencies), while node features evolve over time reflecting market conditions.

\subsection{Cascade Prediction Task}

A \emph{liquidation cascade} is a systemic event where protocol failures propagate through the composability graph, causing correlated TVL losses, price drops, and liquidation spikes across multiple protocols.

\textbf{Definition 1} (Cascade Event). A cascade event $c = (t_s, t_p, t_e, s, \ell)$ is characterized by a start time $t_s$, peak time $t_p$, end time $t_e$, severity $s \in \{\text{catastrophic}, \text{severe}, \text{moderate}\}$, and aggregate TVL loss percentage $\ell \in (0, 1]$.

For each prediction horizon $h \in \mathcal{H} = \{1, 3, 7, 30\}$ days, we define the binary cascade label:
\begin{equation}
    y_t^{(h)} = \begin{cases}
        1 & \text{if } \exists \text{ cascade active within } [t, t+h] \\
        0 & \text{otherwise}
    \end{cases}
\end{equation}

The task is to learn a function $f_\theta: (\mathcal{G}_t, \mathbf{M}_{t-1}) \rightarrow \hat{y}_t^{(h)} \in [0,1]$ for each $h \in \mathcal{H}$, where $\mathbf{M}_{t-1}$ is a memory state encoding the temporal history of the graph.

\subsection{Composability Graph Edge Types}

We define six edge types capturing distinct contagion channels:

\begin{enumerate}
    \item \textbf{Shared Collateral} ($e_{\text{sc}}$): Protocols accepting the same collateral assets (WETH, WBTC, USDC, USDT, DAI, stETH). A price drop in a shared collateral triggers simultaneous liquidations across all protocols accepting it.

    \item \textbf{Liquidity Flow} ($e_{\text{lf}}$): Capital flow pathways between protocols (e.g., Aave $\rightarrow$ Uniswap for liquidation routes, Curve $\rightarrow$ Convex for LP staking).

    \item \textbf{Oracle Dependency} ($e_{\text{od}}$): Protocols relying on the same price oracle (Chainlink, Pyth, TWAP). Oracle manipulation or failure simultaneously affects all dependent protocols.

    \item \textbf{Governance Overlap} ($e_{\text{go}}$): Shared governance token exposure (CRV, CVX, AAVE, COMP, UNI, LDO), creating governance attack vectors and correlated incentive structures.

    \item \textbf{Price Correlation} ($e_{\text{pc}}$): Protocols whose token prices exhibit correlation above 0.7, indicating exposure to common market factors.

    \item \textbf{Liquidation Pathway} ($e_{\text{lp}}$): Directed edges from lending protocols to DEXes representing the path that liquidated collateral follows when sold on-market.
\end{enumerate}


% ============================================================
% IV. Methodology
% ============================================================
\section{Methodology}

Our framework consists of four components: (1) feature engineering, (2) composability graph construction, (3) temporal graph network architecture, and (4) multi-horizon training with focal loss.

\subsection{Feature Engineering}

For each protocol $v_i$ at timestep $t$, we compute a feature vector $\mathbf{x}_{t,i} \in \mathbb{R}^{46}$ organized into six semantic groups:

\textbf{TVL Features} ($F = 8$): Log-TVL, 1/7/30-day percentage changes, drawdown from peak, cross-sectional rank, 90-day z-score, and 30-day moving average ratio. These capture both the level and dynamics of protocol capitalization from real DeFiLlama data.

\textbf{Price Features} ($F = 9$): Log-price, 1/7/30-day returns, 7/30-day volatility, volume proxy, market cap proxy, and price drawdown. Price dynamics are derived from TVL return decomposition, reflecting the joint influence of asset prices and capital flows on protocol TVL.

\textbf{Liquidity Features} ($F = 8$): Utilization ratio, borrow/supply rates, pool depth, effective liquidity, rate spread. Lending market stress indicators are derived from TVL drawdown dynamics, capturing the empirical relationship between capital outflows and utilization spikes.

\textbf{Network Features} ($F = 6$): Degree, betweenness, eigenvector centrality, clustering coefficient, PageRank, and shared collateral count, computed from the static composability graph.

\textbf{Macro Features} ($F = 9$): Federal Funds Rate, Treasury yields (10Y, 2Y, 3M), CPI, M2 money supply, VIX, DXY, S\&P 500. Macroeconomic indicators are reconstructed from publicly reported values using piecewise-linear interpolation anchored at known data points.

\textbf{Temporal Features} ($F = 6$): Cyclical day-of-week and month encodings via sine/cosine transforms, plus a linear trend indicator.

All features are standardized to zero mean and unit variance \emph{using only training-set statistics} to prevent data leakage, with values clipped to $[-10, 10]$.

\subsubsection{Temporal Feature Augmentation}

To provide the model with explicit temporal context beyond what is stored in its memory module, we augment each node's feature vector with rolling window statistics over $W = 10$ timesteps:

\begin{equation}
    \tilde{\mathbf{x}}_{t,i} = \left[\mathbf{x}_{t,i} \;\|\; \bar{\mathbf{x}}_{t,i}^{(W)} \;\|\; \sigma\mathbf{x}_{t,i}^{(W)} \;\|\; \Delta\mathbf{x}_{t,i}^{(W)}\right]
\end{equation}

where $\bar{\mathbf{x}}_{t,i}^{(W)}$ is the $W$-day rolling mean, $\sigma\mathbf{x}_{t,i}^{(W)}$ is the rolling standard deviation, and $\Delta\mathbf{x}_{t,i}^{(W)} = \mathbf{x}_{t,i} - \bar{\mathbf{x}}_{t,i}^{(W)}$ is the deviation from the rolling mean. This augmentation expands the feature dimension from 46 to 184 while providing the model with ready-to-use temporal gradient information.

\subsection{Composability Graph Construction}

The static composability graph is constructed from domain knowledge of DeFi protocol interactions, organized into six edge types as defined in Section~III-C.

For shared collateral edges, we enumerate the major collateral assets (WETH, WBTC, USDC, USDT, DAI, stETH, wstETH) and create bidirectional edges between all protocols accepting each asset. Oracle dependency edges connect protocols sharing common price feed sources. Liquidity flow edges capture empirically observed capital pathways. Governance overlap edges connect protocols with shared governance token exposure. Liquidation pathway edges connect lending protocols to DEXes where liquidated collateral is sold.

The heterogeneous edge index $\mathcal{E}_k$ for edge type $k$ is represented as a sparse tensor of shape $(2, |\mathcal{E}_k|)$. For the homogeneous baseline, we merge all edge types into a single adjacency matrix $\mathbf{A} \in \mathbb{R}^{N \times N}$ with edge-type-specific weights.

\subsection{Temporal Graph Network Architecture}

Our TGN architecture processes each timestep $t$ through five modules:

\subsubsection{Input Projection}
The augmented node features $\tilde{\mathbf{X}}_t \in \mathbb{R}^{N \times 184}$ are projected to a $d$-dimensional embedding space:
\begin{equation}
    \mathbf{H}_t^{(0)} = \text{GELU}\left(\text{LayerNorm}\left(\mathbf{W}_{\text{in}} \tilde{\mathbf{X}}_t + \mathbf{b}_{\text{in}}\right)\right)
\end{equation}
where $d = 128$ is the embedding dimension.

\subsubsection{Memory Module}
Each node $v_i$ maintains a memory state $\mathbf{m}_{t,i} \in \mathbb{R}^{d_m}$ ($d_m = 128$) that is updated after each timestep using a GRU cell:
\begin{equation}
    \mathbf{m}_{t,i} = \text{GRU}\left(\mathbf{h}_{t,i}, \mathbf{m}_{t-1,i}\right)
\end{equation}
The memory captures the temporal evolution of each protocol's state, enabling the model to detect gradual pre-cascade deterioration patterns (e.g., slowly declining TVL, rising utilization rates) that precede sudden cascade events.

The memory and current embeddings are fused via:
\begin{equation}
    \mathbf{z}_{t,i} = \text{GELU}\left(\text{LayerNorm}\left(\mathbf{W}_f [\mathbf{h}_{t,i}^{(0)} \| \mathbf{m}_{t-1,i}] + \mathbf{b}_f\right)\right)
\end{equation}

\subsubsection{Heterogeneous Graph Attention}
For each edge type $k \in \{1, \ldots, K\}$, we apply a multi-head attention message passing layer:
\begin{equation}
    \alpha_{ij}^{(k,l)} = \frac{\exp\left(\text{LeakyReLU}\left(\mathbf{a}_s^{(l)\top}\mathbf{z}_j + \mathbf{a}_d^{(l)\top}\mathbf{z}_i + \mathbf{a}_e^{(l)\top}\mathbf{e}_{ij}^{(k)}\right)\right)}{\sum_{j' \in \mathcal{N}_i^{(k)}} \exp(\cdot)}
\end{equation}
where $l \in \{1, \ldots, L\}$ indexes attention heads ($L = 4$), $\mathbf{e}_{ij}^{(k)}$ are edge features, and $\mathcal{N}_i^{(k)}$ is node $i$'s neighborhood under edge type $k$.

Messages from all edge types are aggregated:
\begin{equation}
    \mathbf{z}_{t,i}^{(\ell+1)} = \text{LayerNorm}\left(\mathbf{W}_{\text{agg}} \left[\bigoplus_{k=1}^K \sum_{j \in \mathcal{N}_i^{(k)}} \alpha_{ij}^{(k)} \mathbf{W}_k \mathbf{z}_j^{(\ell)}\right]\right)
\end{equation}
where $\bigoplus$ denotes concatenation and $\ell$ indexes GNN layers ($\ell \in \{1, 2\}$).

\subsubsection{Graph-Level Readout}
We obtain a graph-level representation via attention-weighted pooling:
\begin{equation}
    \mathbf{g}_t = \sum_{i=1}^N \text{softmax}\left(\mathbf{w}_a^\top \mathbf{z}_{t,i}^{(L)}\right) \cdot \mathbf{z}_{t,i}^{(L)}
\end{equation}
Combined with max-pooling for extreme value detection:
\begin{equation}
    \mathbf{r}_t = [\mathbf{g}_t \| \max_i \mathbf{z}_{t,i}^{(L)}] \in \mathbb{R}^{2d}
\end{equation}

\subsubsection{Multi-Horizon Prediction}
Separate prediction heads for each horizon $h \in \mathcal{H}$ produce cascade logits:
\begin{equation}
    \hat{y}_t^{(h)} = \sigma\left(\text{MLP}_h(\mathbf{r}_t)\right)
\end{equation}
where each MLP has two hidden layers ($128 \rightarrow 64 \rightarrow 1$) with GELU activations and dropout ($p = 0.15$).

\subsection{Training}

\subsubsection{Focal Loss}
Cascade events are rare ($\sim$5--8\% of timesteps), creating severe class imbalance. We use focal loss~\cite{lin2017focal} with class weighting:
\begin{equation}
    \mathcal{L}_{\text{focal}} = -\alpha_t (1 - p_t)^\gamma \cdot w_t \cdot \log(p_t)
\end{equation}
where $\gamma = 2.0$ is the focusing parameter, $\alpha = 0.75$, $w_t = 10.0$ for positive samples, and $p_t = \sigma(\hat{y})y + (1 - \sigma(\hat{y}))(1-y)$ is the probability assigned to the correct class.

The total loss combines cascade prediction and severity regression:
\begin{equation}
    \mathcal{L} = \sum_{h \in \mathcal{H}} \mathcal{L}_{\text{focal}}^{(h)} + \lambda \cdot \text{MSE}(\hat{s}_t, s_t)
\end{equation}
where $\lambda = 0.3$ and $s_t \in [0, 1]$ is the continuous risk severity score.

\subsubsection{Optimization}
We use AdamW~\cite{loshchilov2019adamw} with learning rate $3 \times 10^{-4}$, weight decay $10^{-4}$, and cosine annealing schedule (minimum LR $10^{-6}$). Gradient norms are clipped to 1.0. Training runs for up to 100 epochs with early stopping (patience = 20) based on validation loss.

\subsubsection{Windowed Truncated BPTT}
Memory states are trained using windowed truncated backpropagation through time (TBPTT) with a window size of 10 timesteps. At each window boundary, gradients are accumulated and applied, and memory states are detached to prevent prohibitively expensive full-sequence backpropagation. Memory persists across epochs, allowing the GRU to build long-term protocol state representations.

\subsubsection{Temporal Data Splits}
We use a chronological split to prevent data leakage, with boundaries chosen to ensure cascade events in each partition:
\begin{itemize}
    \item \textbf{Training} (June 2021--October 2022, 518 days): Covers Terra/Luna collapse and 3AC/Celsius contagion.
    \item \textbf{Validation} (November 2022--March 2023, 151 days): Covers FTX collapse, USDC depeg/SVB crisis, and Euler hack.
    \item \textbf{Test} (April 2023--February 2026, 1,065 days): Covers Curve exploit and statistically detected anomaly events.
\end{itemize}
Feature normalization statistics are computed exclusively from training data.


% ============================================================
% V. Experimental Setup
% ============================================================
\section{Experimental Setup}

\subsection{Dataset}

We use real-world TVL data sourced from DeFiLlama~\cite{defillama} covering 15 major Ethereum DeFi protocols from June 2021 to February 2026 (1,734 daily snapshots). The protocols span five functional categories: lending (Aave v2/v3, Compound v2/v3, MakerDAO, Morpho), DEX (Uniswap v2/v3, Curve, Balancer), liquid staking (Lido, Rocket Pool), yield aggregation (Yearn, Convex), and stablecoin (Frax).

Table~\ref{tab:events} lists the cascade events used for labeling. Six events are verified historical crises with known dates and severity. An additional seven events are detected through statistical anomaly analysis on aggregate TVL (z-score $< -4.5$ or drawdown $> 20\%$, minimum 5-day duration), capturing significant market stress periods not associated with named events.

\begin{table}[!t]
\caption{Cascade Events Used for Labeling}
\label{tab:events}
\centering
\footnotesize
\begin{tabular}{llcc}
\toprule
\textbf{Event} & \textbf{Date} & \textbf{Severity} & \textbf{TVL Loss} \\
\midrule
Terra/Luna & May 2022 & Catastrophic & 45\% \\
3AC/Celsius & Jun 2022 & Severe & 30\% \\
FTX Collapse & Nov 2022 & Severe & 25\% \\
USDC Depeg/SVB & Mar 2023 & Moderate & 12\% \\
Euler Hack & Mar 2023 & Moderate & 8\% \\
Curve Exploit & Jul 2023 & Moderate & 10\% \\
\midrule
\multicolumn{4}{l}{\emph{+ 7 statistically detected anomaly events}} \\
\bottomrule
\end{tabular}
\end{table}

\subsubsection{Feature Sources}

TVL data---the primary predictive signal---is sourced directly from DeFiLlama and used without modification. Supplementary features are derived transparently:
\begin{itemize}
    \item \emph{Price features}: Decomposed from TVL dynamics (TVL $=$ quantity $\times$ price), using a 60/40 blend of market-wide and protocol-specific TVL returns as price proxy.
    \item \emph{Lending features}: Utilization rates derived from TVL drawdown dynamics (capital outflows reduce available supply, increasing utilization), with kinked rate curves.
    \item \emph{Macro features}: Reconstructed from publicly reported values (Federal Funds Rate, CPI, VIX, etc.) using piecewise-linear interpolation at known data points.
    \item \emph{On-chain features}: TVL velocity (daily absolute change magnitude) as transaction activity proxy.
\end{itemize}

The resulting dataset has cascade positive rates of approximately 5--8\% at the 1-day horizon, representing a realistic rare-event prediction setting.

\subsection{Baselines}

We compare against five baselines, each representing a distinct modeling paradigm:

\textbf{Static GNN (GAT).} A graph attention network~\cite{velickovic2018gat} operating on the same composability graph but without temporal components (no memory, no time encoding). Uses the same augmented features as TGN for fair comparison. This baseline isolates the contribution of temporal modeling.

\textbf{LSTM.} A two-layer LSTM~\cite{hochreiter1997lstm} processing 30-day sequences of flattened protocol features. This baseline captures temporal patterns without graph structure, isolating the contribution of graph topology.

\textbf{XGBoost.} Gradient-boosted trees~\cite{chen2016xgboost} operating on protocol-aggregated features (cross-protocol mean, std, min, max) with temporal rolling statistics. This non-deep-learning baseline tests whether the patterns are learnable without representation learning.

\textbf{SIR Contagion.} A Susceptible-Infected-Recovered compartmental model~\cite{hurd2016contagion} simulating cascade propagation on the adjacency matrix. The initial infected state is derived from TVL-based stress indicators. This physics-based baseline tests whether domain-specific epidemiological models suffice.

\textbf{Network Centrality.} Centrality metrics (degree, betweenness, eigenvector, PageRank, clustering coefficient) combined with feature aggregates, fed to logistic regression with balanced class weights. This baseline tests whether simple network analysis captures cascade risk.

\subsection{Evaluation Metrics}

We evaluate using AUROC (primary metric, threshold-independent) and AUPRC (sensitive to class imbalance). For threshold-dependent metrics (F1, precision, recall, specificity), we select the optimal threshold via Youden's J statistic ($J = \text{TPR} - \text{FPR}$), avoiding the common failure mode of fixed 0.5 thresholds on imbalanced data.

Additional metrics include Matthews Correlation Coefficient (MCC), Brier score (calibration), and confusion matrix statistics.

\subsection{Statistical Testing}

We validate performance differences using:
\begin{itemize}
    \item \textbf{Diebold-Mariano test}~\cite{diebold2002comparing}: Compares forecast accuracy between model pairs using squared error differences, with HAC standard errors.
    \item \textbf{McNemar's test}: Compares classification accuracy using contingency tables of correct/incorrect predictions.
    \item \textbf{Bootstrap confidence intervals}: 3{,}000 iterations with Bonferroni correction for multiple comparisons.
\end{itemize}


% ============================================================
% VI. Results
% ============================================================
\section{Results}

\subsection{Main Results}

Table~\ref{tab:main_results} reports AUROC and AUPRC across all models and prediction horizons on the held-out test set (April 2023--February 2026).

\begin{table*}[!t]
\caption{Model Comparison: AUROC and AUPRC Across Prediction Horizons}
\label{tab:main_results}
\centering
\begin{tabular}{l|cc|cc|cc|cc}
\toprule
& \multicolumn{2}{c|}{\textbf{1-day}} & \multicolumn{2}{c|}{\textbf{3-day}} & \multicolumn{2}{c|}{\textbf{7-day}} & \multicolumn{2}{c}{\textbf{30-day}} \\
\textbf{Model} & AUROC & AUPRC & AUROC & AUPRC & AUROC & AUPRC & AUROC & AUPRC \\
\midrule
\textbf{TGN (Ours)} & \TODO{X.XXX} & \TODO{X.XXX} & \TODO{X.XXX} & \TODO{X.XXX} & \TODO{X.XXX} & \TODO{X.XXX} & \TODO{X.XXX} & \TODO{X.XXX} \\
Static GNN & \TODO{X.XXX} & \TODO{X.XXX} & \TODO{X.XXX} & \TODO{X.XXX} & \TODO{X.XXX} & \TODO{X.XXX} & \TODO{X.XXX} & \TODO{X.XXX} \\
LSTM & \TODO{X.XXX} & \TODO{X.XXX} & \TODO{X.XXX} & \TODO{X.XXX} & \TODO{X.XXX} & \TODO{X.XXX} & \TODO{X.XXX} & \TODO{X.XXX} \\
XGBoost & \TODO{X.XXX} & \TODO{X.XXX} & \TODO{X.XXX} & \TODO{X.XXX} & \TODO{X.XXX} & \TODO{X.XXX} & \TODO{X.XXX} & \TODO{X.XXX} \\
SIR Contagion & \TODO{X.XXX} & \TODO{X.XXX} & \TODO{X.XXX} & \TODO{X.XXX} & \TODO{X.XXX} & \TODO{X.XXX} & \TODO{X.XXX} & \TODO{X.XXX} \\
Centrality & \TODO{X.XXX} & \TODO{X.XXX} & \TODO{X.XXX} & \TODO{X.XXX} & \TODO{X.XXX} & \TODO{X.XXX} & \TODO{X.XXX} & \TODO{X.XXX} \\
\bottomrule
\end{tabular}
\end{table*}

\TODO{Narrative: Describe main findings. Which models perform best at each horizon? Does graph structure help? Does temporal modeling help? How do short vs long horizons compare?}

\subsection{Statistical Significance}

Table~\ref{tab:statistical} reports Diebold-Mariano test p-values comparing TGN against each baseline.

\begin{table}[!t]
\caption{Statistical Significance: Diebold-Mariano Test p-values (TGN vs. Baselines)}
\label{tab:statistical}
\centering
\footnotesize
\begin{tabular}{lcccc}
\toprule
\textbf{TGN vs.} & \textbf{1d} & \textbf{3d} & \textbf{7d} & \textbf{30d} \\
\midrule
Static GNN & \TODO{p} & \TODO{p} & \TODO{p} & \TODO{p} \\
LSTM & \TODO{p} & \TODO{p} & \TODO{p} & \TODO{p} \\
XGBoost & \TODO{p} & \TODO{p} & \TODO{p} & \TODO{p} \\
SIR & \TODO{p} & \TODO{p} & \TODO{p} & \TODO{p} \\
Centrality & \TODO{p} & \TODO{p} & \TODO{p} & \TODO{p} \\
\bottomrule
\end{tabular}
\end{table}

\TODO{Bootstrap CIs for TGN AUROC per horizon. Significance narrative.}

\subsection{Ablation Studies}

\subsubsection{Feature Group Ablation}

Table~\ref{tab:ablation_feat} shows the change in TGN AUROC when each feature group is zeroed out.

\begin{table}[!t]
\caption{Feature Group Ablation: $\Delta$AUROC when Removing Each Group}
\label{tab:ablation_feat}
\centering
\footnotesize
\begin{tabular}{lcccc}
\toprule
\textbf{Removed Group} & \textbf{1d} & \textbf{3d} & \textbf{7d} & \textbf{30d} \\
\midrule
TVL Features & \TODO{$\Delta$} & \TODO{$\Delta$} & \TODO{$\Delta$} & \TODO{$\Delta$} \\
Price Features & \TODO{$\Delta$} & \TODO{$\Delta$} & \TODO{$\Delta$} & \TODO{$\Delta$} \\
Liquidity Features & \TODO{$\Delta$} & \TODO{$\Delta$} & \TODO{$\Delta$} & \TODO{$\Delta$} \\
Network Features & \TODO{$\Delta$} & \TODO{$\Delta$} & \TODO{$\Delta$} & \TODO{$\Delta$} \\
Macro Features & \TODO{$\Delta$} & \TODO{$\Delta$} & \TODO{$\Delta$} & \TODO{$\Delta$} \\
Temporal Features & \TODO{$\Delta$} & \TODO{$\Delta$} & \TODO{$\Delta$} & \TODO{$\Delta$} \\
\bottomrule
\end{tabular}
\end{table}

\subsubsection{Edge Type Ablation}

Table~\ref{tab:ablation_edge} shows the impact of removing each edge type from the composability graph.

\begin{table}[!t]
\caption{Edge Type Ablation: $\Delta$AUROC when Removing Each Edge Type}
\label{tab:ablation_edge}
\centering
\footnotesize
\begin{tabular}{lcccc}
\toprule
\textbf{Removed Edge Type} & \textbf{1d} & \textbf{3d} & \textbf{7d} & \textbf{30d} \\
\midrule
Shared Collateral & \TODO{$\Delta$} & \TODO{$\Delta$} & \TODO{$\Delta$} & \TODO{$\Delta$} \\
Liquidity Flow & \TODO{$\Delta$} & \TODO{$\Delta$} & \TODO{$\Delta$} & \TODO{$\Delta$} \\
Oracle Dependency & \TODO{$\Delta$} & \TODO{$\Delta$} & \TODO{$\Delta$} & \TODO{$\Delta$} \\
Governance Overlap & \TODO{$\Delta$} & \TODO{$\Delta$} & \TODO{$\Delta$} & \TODO{$\Delta$} \\
Price Correlation & \TODO{$\Delta$} & \TODO{$\Delta$} & \TODO{$\Delta$} & \TODO{$\Delta$} \\
Liquidation Pathway & \TODO{$\Delta$} & \TODO{$\Delta$} & \TODO{$\Delta$} & \TODO{$\Delta$} \\
\bottomrule
\end{tabular}
\end{table}

\subsubsection{Component Ablation}
Removing the memory module results in a $\Delta$AUROC of \TODO{$\Delta$} averaged across horizons, confirming that temporal memory tracking is critical for cascade prediction.

\TODO{Narrative: Which feature groups matter most? Which edges? Memory contribution.}


% ============================================================
% VII. Discussion
% ============================================================
\section{Discussion}

\subsection{DeFi as a Computational Social System}

Our results highlight that DeFi cascade prediction benefits from modeling the ecosystem as a computational social system.
The composability graph---with its six edge types capturing shared collateral, liquidity flows, oracle dependencies, governance overlaps, price correlations, and liquidation pathways---is fundamentally a reflection of human design decisions and community governance.
Each edge represents a dependency that was created through DAO votes, protocol integrations approved by governance, or emergent patterns of collective capital allocation.

The edge type ablation results reveal which social-technical dependency channels contribute most to cascade prediction.
\TODO{Discuss which edges matter most and why, relating to the social governance structure.}

Governance overlap edges are particularly noteworthy from a computational social systems perspective: they capture the phenomenon where overlapping community membership across protocols creates correlated risk-taking behavior.
When a governance token (e.g., CRV) experiences a crisis, it simultaneously affects the governance capacity and community confidence of multiple protocols, illustrating how social dynamics drive technical contagion.

\subsection{Why Graph Structure Matters}

The composability graph provides critical information for cascade prediction because contagion propagates through specific dependency channels.
\TODO{Discuss TGN vs LSTM comparison showing graph advantage. Discuss edge ablation showing which contagion channels are most important.}

\subsection{Temporal Memory for Pre-Cascade Detection}

The GRU memory module enables TGN to track gradual protocol deterioration patterns---slowly declining TVL, rising utilization, increasing rate spreads---that precede sudden cascade events.
\TODO{Discuss TGN vs Static GNN showing temporal advantage. Memory ablation results.}

\subsection{Multi-Horizon Prediction Patterns}

Short-horizon predictions (1-day, 3-day) benefit from stronger immediate pre-cascade signals, while the 30-day horizon requires detecting early warning signs amid normal market noise.
\TODO{Discuss horizon-specific performance patterns.}

\subsection{Practical Implications}

Our framework has direct applications in DeFi risk management:
\begin{itemize}
    \item \textbf{Operational monitoring}: 1-day and 3-day predictions can trigger automated position adjustments and liquidity management.
    \item \textbf{Strategic planning}: 7-day and 30-day predictions inform portfolio rebalancing and governance proposals.
    \item \textbf{Protocol governance}: Risk dashboards can visualize per-protocol cascade vulnerability, enabling DAO members to make informed governance decisions---a direct application within the computational social systems paradigm.
\end{itemize}

\subsection{Limitations}

\textbf{Supplementary feature derivation.} While TVL data is directly sourced from DeFiLlama, supplementary features (prices, lending metrics, macro) are derived from TVL dynamics or reconstructed from known values rather than collected from primary sources. Future work should incorporate direct on-chain data, oracle price feeds, and lending protocol APIs for fully granular feature sets.

\textbf{Static graph.} The composability graph is currently static, but real DeFi dependencies evolve as protocols upgrade, new integrations launch, and liquidity migrates. Extending to dynamic graph construction that reflects governance-driven changes in protocol relationships is an important direction.

\textbf{Protocol coverage.} Our 15-protocol graph covers major protocols but excludes the long tail of smaller protocols that may serve as cascade amplifiers or early warning canaries.

\textbf{Temporal resolution.} Daily snapshots may miss intra-day dynamics during fast-moving cascade events. Higher-frequency data (hourly or block-level) would improve short-horizon predictions.


% ============================================================
% VIII. Conclusion
% ============================================================
\section{Conclusion}

We presented a Temporal Graph Network framework for predicting DeFi liquidation cascades, framing the DeFi ecosystem as a computational social system where protocol communities interact through technically and socially mediated dependency channels.
By modeling the ecosystem as a heterogeneous temporal graph with six edge types capturing distinct contagion channels, and combining per-node GRU memory with multi-head graph attention, our model jointly captures the spatial structure and temporal dynamics of cascade propagation.
Evaluated on real-world TVL data spanning 1,734 days across 15 major DeFi protocols, encompassing six verified historical cascade events and statistically detected anomalies, our approach demonstrates the value of jointly modeling graph topology and temporal dynamics for cascade prediction.
Ablation studies confirm the complementary contributions of different composability graph edge types, temporal memory, and multi-scale feature engineering.

Our work contributes to the computational social systems literature by demonstrating how graph-based machine learning can model cascade contagion in decentralized financial communities, where social governance structures, collective capital allocation decisions, and technical protocol interdependencies jointly determine systemic risk dynamics.

Future work will extend this framework to fully on-chain data, dynamic graph construction reflecting governance-driven dependency changes, higher-frequency temporal resolution, and integration with automated risk management systems for DAO governance.


% ============================================================
% References
% ============================================================
\bibliographystyle{IEEEtran}
\bibliography{references}

\end{document}
